\section{Introducción}
El problema de las $n$ reinas pregunta cuántas maneras hay de colocar $n$ reinas en un
tablero de $n\times n$ sin que dos se ataquen. Denotamos ese número por $Q(n)$ y contamos
por separado las soluciones relacionadas por una simetría. Leer las columnas de sus reinas
por orden de filas identifica cada solución con una permutación. Las restricciones diagonales
determinan cuáles de las $n!$ permutaciones sirven.

Simkin~\cite[Teorema~1.1]{simkin} obtuvo el orden principal de $\log Q(n)$ mediante
un problema variacional. Los resultados de Luria y Simkin~\cite{ls} y de Bowtell y
Keevash~\cite{bk} estudian el conteo por métodos combinatorios. Nobel, Agrawal y Boyd~\cite{nab}
acotaron la constante de primer orden. Para determinar el factor multiplicativo que queda en
el conteo después de extraer el término principal hace falta una estimación más precisa.

Sea $Q_{\rm pred}(n)$ el predictor canónico de~\cite{pred} y pongamos
$R_Q(n)=\log Q(n)-\log Q_{\rm pred}(n)$. El siguiente resultado establece que
$R_Q(n)\to0$, la forma cualitativa de la Conjetura~10.1 de~\cite{pred}.
\begin{teorema}[Término constante del conteo]\label{princ}
En el tablero ordinario, cuando $n\to\infty$,
\[
 \log Q(n)-\log n!-n\Phi\longrightarrow
 \tfrac12\log2-\tfrac12-\tfrac32 I_2-\tfrac12\log{\det}_2(I+J)=C_{\rm pred}.
\]
\end{teorema}
Aquí $\nu_1$ es el perfil óptimo de masa uno de~\cite{perfil}, $\Phi=s_1+1$ con $s_1$ su entropía, e
$I_2=\int_{[0,1]^2}\nu_1^2$. El operador $J$ es el de (D3.6) de~\cite{pred}, y su determinante
regularizado se describe en (D4.1) de ese trabajo. Un error que tiende a cero en el logaritmo
equivale a un cociente de conteos que tiende
a uno. La estimación cuantitativa $|R_Q(n)|\le C/n$ de la misma conjetura sigue abierta.

En la Sección~9 de~\cite{pred} se estima $C_{\rm pred}\approx0{,}9401$, y la evaluación directa de
la fórmula da $0{,}940144$; ninguna de las dos cifras tiene cota de error certificada. Con la
referencia $\Phi_0=1-1{,}944001$ de esa sección, los conteos exactos para $16\le n\le27$ dan
$\log Q(n)-\log n!-n\Phi_0$ entre $0{,}913$ y $0{,}941$ \cite[Tabla~2]{pred}, y $R_Q(n)$ vale
$0{,}0208$ en $n=20$ y $0{,}0119$ en $n=27$. Estos valores son compatibles con el
Teorema~\ref{princ}, pero ningún cálculo en valores finitos de $n$ lo comprueba, porque el
teorema afirma un límite sin velocidad de convergencia.

La Sección~10 de~\cite{pred} identifica tres cuestiones. La primera es la regularidad de
los tableros residuales a lo largo del proceso aleatorio. La segunda es el efecto de las
últimas filas, donde quedan pocas opciones de inserción. La tercera es cómo reunir los errores
de las sucesivas inserciones al comparar ambos conteos.

\paragraph{Esquema de la demostración.}
Escribimos $g(B)=\log Q(B)-L(B)$, donde $L(B)$ es el predictor de un tablero residual $B$
(Sección~\ref{leyes}). Hay que probar que $g(B_0)\to0$ en el tablero inicial $B_0$. Comparamos dos
leyes sobre las historias que insertan una reina por fila, en un orden uniforme de las filas. La ley
uniforme $\mu$ revela una solución uniforme. La ley adaptativa $\nu$ elige cada casilla $e$ con
probabilidad proporcional a $\exp L(B+e)$, la predicción del tablero que resulta. La
Sección~\ref{historias} expresa $g(B_0)$ con ambas leyes en una historia detenida: una identidad con
la entropía relativa de $\mu$ respecto de $\nu$ y una cota inferior bajo $\nu$. El término acumulado
de esas fórmulas es pequeño por la consistencia de~\cite{pred}.

Fijamos $\delta>0$ pequeño y un entero $K$ grande, y ponemos $M=\lceil\delta n\rceil$. El orden en
que se eligen las constantes está en la Sección~\ref{meso:constantes}. Las Secciones~\ref{macro}
y~\ref{meso} prueban que, bajo $\nu$, los tableros de la historia son regulares
(Definición~\ref{def:regular}) hasta que quedan $K$ filas vacías, salvo probabilidad
$O(e^{-cK})+o(1)$: primero mientras quedan al menos $M$ filas y después entre $M$ y $K$. La
Sección~\ref{costes} acota por arriba $g$ en el primer tablero que deja de ser regular. La
Sección~\ref{terminal} prueba que, en el tablero de $K$ filas, el conteo es un permanente fuera de
un suceso de probabilidad $O(K^4/n)$, y que allí $g$ tiende a cero al crecer $K$. En la
Sección~\ref{balance} se hace primero $n\to\infty$ con $\delta$ y $K$ fijos y después $K\to\infty$.

Los Teoremas~1.1--1.3, la Proposición~5.2 y el Teorema~7.7 de~\cite{perfil} proporcionan el
perfil continuo y sus cotas. La comparación entrópica del Teorema~3.6 de~\cite{ultimas}
motiva el argumento de la Sección~\ref{costes}, donde verificamos las estimaciones sobre el
soporte de cada residual. Usamos los Teoremas~7.1 y~8.1 y el Corolario~7.2 de~\cite{pred}
para las inserciones, y sus Teoremas~1.1 y~5.7 para la constante. La única estimación
asintótica de permanentes que usamos es la de McCullagh~\cite{mccullagh}, en la forma uniforme
probada por Divoux y Wang~\cite[Proposición~3.1 y Observación~2.1]{dw}. En la
Sección~\ref{permanente} verificamos sus hipótesis tras escalar la matriz del residual.
